% Discussion: the answer to the Introduction's question, the accounts weighed (Table 2), what the
% questions asked, the law as parallels and divergences, and the limits. Planned in
% notes/discussion_plan.md (21 October 2026); every number through a macro; statuses follow
% notes/argument.md and notes/results_ledger.md. American spelling.

Asked about an image that a label tied to an AI system, people gave it less credit; asked about a person's own rework of that image, they gave its AI origin little weight. The AI label lowered ratings of the same images, more for effort in Study~1, and a stated 100 hours did not close the gap. Once Alex had edited each image by hand, ratings followed how much he had changed it, and the label carried much less into them than into ratings of the originals. That is a reading across studies, not a test, and the carry-over was not zero: it was small, uneven and detectable at slight edits in Study~5. These ratings do not support the original-sin hypothesis in its strong form. An AI label behaved more like information about who did what than like a stain on the object.

% Table 2: the accounts the Introduction names, what each predicts, what the five studies show and
% how the evidence reads. Written by hand from notes/argument.md (claim 15) and the Results; every
% number is a macro from generated/numbers.tex. Statuses follow notes/results_ledger.md.
\begin{table}[tbp]
\centering
\caption{\textbf{Accounts of the AI label's effect, and what the five studies show.} The penalty is the mean rating under the Human label minus the mean under the AI label. Comparisons between originals and second works are readings across studies, not tests. Scoped credit was stated after the results were known.}
\label{tab:accounts}
\small
\setlength{\tabcolsep}{4pt}
\begin{tabular}{@{}p{70pt}p{92pt}p{\dimexpr\textwidth-262pt-6\tabcolsep\relax}p{100pt}@{}}
\toprule
Account & What it predicts & What the studies show & Reading\tabularnewline
\midrule
\raggedright \textbf{Effort signal}: the label only signals little work & \raggedright A stated 100 hours closes the gap on originals & \raggedright The registered decision on both ratings: the gap persists; its 90\% CIs at 100 hours reach below $\pm$\SourceBound{} (Study~3) & \raggedright Not the whole story; a partial role, above all on effort, is not excluded\tabularnewline
\addlinespace[3pt]
\raggedright \textbf{Taint}: an AI origin passes to whatever is built from it & \raggedright The originals' penalty, at full strength, at every level of edit & \raggedright Every registered interval for the creativity penalty on second works excludes Study~1's \SOneHOneCreatDiffPrecise{} points & \raggedright Not supported at full strength\tabularnewline
\addlinespace[3pt]
\raggedright \textbf{Dilution}: the penalty shrinks as the edit grows & \raggedright Largest at None, falling toward Dramatic & \raggedright Study~2's registered test was not significant and ran the other way; no fall in the pooled model (exploratory); the floor at None weakens both & \raggedright Not supported\tabularnewline
\addlinespace[3pt]
\raggedright \textbf{Severance}: the label does not reach second works & \raggedright No penalty at any level & \raggedright A penalty at Slight on all three ratings in Study~5 (registered secondary, uncorrected); BIC favored no penalty in the pooled model (exploratory) & \raggedright Not supported in its strict form, but close to the data\tabularnewline
\addlinespace[3pt]
\raggedright \textbf{Visible traces}: ratings follow what the image shows & \raggedright No label effect on identical images & \raggedright A penalty on identical originals in Studies 1 and 3 (registered confirmatory) & \raggedright Cannot explain the originals' penalty\tabularnewline
\addlinespace[3pt]
\raggedright \textbf{Contagion}: an object carries something of how it was made & \raggedright A penalty on originals that also reaches what is built from them & \raggedright A penalty on originals; a small, uneven carry-over to second works & \raggedright Not separable from scoped credit on originals; a strong version is weakened\tabularnewline
\addlinespace[3pt]
\raggedright \textbf{Scoped credit}: credit follows the act the label describes & \raggedright A penalty on originals; stated time raises credit for the act described; little carry-over to a different act; stated hours of the edit raise credit for the edit & \raggedright The first three held; stated editing hours showed no detectable effect (creativity \SFiveCreatEditTimeContrastEst{} points, 90\% CI \SFiveCreatEditTimeContrastLoNinety{} to \SFiveCreatEditTimeContrastHiNinety; Study~5, registered) & \raggedright Organizes most results; one failed prediction\tabularnewline
\bottomrule
\end{tabular}
\end{table}


Table~\ref{tab:accounts} sets the accounts named in the Introduction against this evidence. The label did not work only by signaling little work, since a stated 100 hours left the gap in place, though a partial role for that signal, above all on effort, is not excluded. Neither taint at full strength nor dilution is supported. Severance comes closest: one standard score for choosing between models favored no penalty on second works at all, and the penalties that stand against it come from uncorrected tests, most of them at slight edits. On the originals the images were identical, so the penalty came from the label, but these data cannot separate the account we favor from contagion, the belief that an object carries something of how it was made. That account, scoped credit, which we stated after seeing the results, holds that people credit a person for what they take that person to have done and read a label as saying which act produced what they see. It predicted a penalty on originals, credit that rises with the time stated for the act the label describes, and little carry-over to a different act, and each held. It also predicted that hours stated for Alex's edit would raise credit for the edit, and they did not (creativity \SFiveCreatEditTimeContrastEst{} points, 90\% CI \SFiveCreatEditTimeContrastLoNinety{} to \SFiveCreatEditTimeContrastHiNinety). A single stated time may be hard to weigh without another to compare it with, or a visible edit may leave a stated time little to add. Both readings came after the fact; stating several editing times to the same people would separate them.

The most important limit on this answer is what the questions asked. The second-work items asked about the new work and the edit, so a rater who answered them literally had little reason to use the label. Scoped credit and a literal reading therefore predict the same small carry-over, and these studies cannot tell them apart. A closer study suggests that the difference matters. When an artist painted the same landscape by hand and viewers were told that an AI tool had supplied the compositions, ratings of the painting as a whole fell on most measures, while the artist's creativity was rated higher.\cite{messer2024cocreating} Together, the two sets of results suggest that an AI origin may count against a whole work but not against a person's contribution when the question isolates that contribution, a hypothesis that asking both questions about the same pairs would test. The legal questions are scoped in just this way: a registration claims only what a person contributed, with the AI material disclaimed.\cite{usco2023guidance}

For the law, the ratings show parallels and divergences, not a match. The label lowered effort more than creativity in Study~1, and \emph{Feist} set labor aside as a ground of originality while keeping a minimal creative spark;\cite{feist1991rural} but Human-label effort explanations named visible detail far more often than time or skill (exploratory), so the effort rating is no measure of labor. Stated hours raised creativity ratings of originals by \SThreeCreatTimeContrastEst{} points, a departure from \emph{Feist}'s rejection of labor, though under the AI label more time also meant more control over the image. On second works, the disclosed AI origin counted for little in ratings of the new material, a parallel to the Office's practice of registering a person's contribution with the AI material disclaimed. It has been argued that the market's penalty on AI-labeled work, with the human-authorship rule, gives makers reasons to hide their use of AI;\cite{notivictor2025hidden} here, disclosing an AI source cost a person's rework little in ratings, though only for students asked about that work. The creativity item did not say whose creativity, so the ratings say nothing about whether prompting alone makes a person an author. We claim relevance only in narrow senses: originality can reach a jury,\cite{eleventhcircuit2025pattern} the line between prompting and authorship is unsettled, and a gap between a rule and the intuitions of the people it governs is worth knowing.

The remaining carry-over has no single cause. It was largest at slight edits in Study~5, mostly among people told of a 10-hour edit; it appeared at dramatic edits in Studies 2 and 5; and in Study~4 it ran the other way, not significantly (Supplementary Note~6). Every registration followed data collection. The participants were \NPeople{} students from one subject pool, \NPeopleRepeat{} of them in two studies; the ten images came from one generator; nobody checked belief in the cover story; the debriefing overstated the law; effort was never rated on second works; some Slight versions were also cropped; and no second work had a source of neutral origin. The study most likely to change our conclusion would ask, about the same pairs, how creative the whole work is and how creative the person's contribution is. Others would vary editing time within each person, describe the AI as a tool rather than the maker, and ask the people who apply the rules, such as examiners and jurors.
